\section{Implementação, Determinismo e Terminação}

\subsection{Arquitetura e Fonte Única de Verdade}

A aplicação foi implementada em \textbf{Python 3} (usando exclusivamente a biblioteca padrão para o motor formal, a CLI e a Interface Web Interativa, e \texttt{pytest} para a suíte automatizada de 92 testes da AV2), mantendo ativo o \texttt{MiniLexer} da AV1 em \texttt{src/} e arquivando os entregáveis legados da AV1 em \texttt{archive/av1/}.

O módulo \texttt{src/av2/spec.py} concentra, em estruturas imutáveis (\texttt{@dataclass(frozen=True, slots=True)}), as especificações canônicas da GLC, do AP e do ALL. A partir dele operam:
\begin{itemize}
    \item \texttt{src/av2/glc.py}: motor de derivação mais à esquerda LL(1);
    \item \texttt{src/av2/ap.py}: simulador formal de DPDA com pilha real (\texttt{list[str]}), validação de $\Sigma_{\text{AP}}$ e rastreamento completo;
    \item \texttt{src/av2/all\_lba.py}: simulador formal de D-LBA com fita fixa \texttt{['<', *entrada, '>']}, asserções de fronteira a cada passo e rastreamento completo;
    \item \texttt{src/av2/pipeline.py}: orquestração das duas etapas (o ALL só executa se o AP aceitar);
    \item \texttt{src/av2/minilang\_source.py}: ponte transparente via \texttt{MiniLexer} que projeta código \texttt{MiniLang} em $\Sigma_{\text{AP}}$ sem checagem sintática/semântica externa, mapeia \texttt{(linha, coluna)} para sublinhado de erro estilo IDE e sintetiza código para os 50 casos oficiais;
    \item \texttt{src/av2/cli.py} e \texttt{src/av2/web.py}: CLI e Web IDE estilo terminal (\texttt{./run}, \texttt{./ap}, \texttt{./all}, \texttt{./glc}, \texttt{./suite}, \texttt{./specs}, \texttt{./man}) com grafos SVG ao vivo, modo \textit{fullscreen}, pilha vertical, fita limitada e derivação;
    \item \texttt{src/av2/jflap\_io.py}: gerador e validador estrutural de \texttt{glc.jff}, \texttt{ap.jff} e \texttt{all.jff} compatíveis com o JFLAP 7.1~\cite{jflap}.
\end{itemize}

Em ambos os autômatos, a decisão resulta \textbf{estritamente} da consulta às tabelas $\delta_{\text{AP}}$ e $\delta_{\text{ALL}}$ e das operações reais de pilha e fita --- sem expressões regulares, contadores externos ou comparações diretas de strings.

\subsection{Determinismo e Ausência de Ramificação}

Tanto o AP quanto o ALL foram projetados como máquinas \textbf{100\% determinísticas}:
\begin{itemize}
    \item No \textbf{AP}, não existe par $(q, a, X)$ com mais de uma transição, e em $q_2$ nenhuma leitura atua sobre o topo $Z$; logo, a única $\varepsilon$-transição $(q_2, \varepsilon, Z) \to (q_3, Z)$ jamais conflita com leituras. O detector nativo \texttt{PDANondeterminismDetector} do JFLAP 7.1 confirmou \textbf{0 estados não-determinísticos} em \texttt{ap.jff}.
    \item No \textbf{ALL}, cada par $(q, X)$ mapeia para no máximo uma tripla $(q', Y, D)$, dispensando \textit{backtracking}: existe exatamente um caminho computacional por entrada e a rejeição ocorre imediatamente na ausência de transição.
\end{itemize}

\subsection{Prova de Terminação}

\begin{itemize}
    \item \textbf{Terminação do AP ($O(n)$):} toda transição de $\delta_{\text{AP}}$ avança a leitura em $+1$, exceto $(q_2, \varepsilon, Z) \to (q_3, Z)$. Como não existe retorno de $\{q_3, \dots, q_8\}$ para $q_2$, essa $\varepsilon$-transição ocorre no máximo uma vez por execução. Para entrada de tamanho $n$, o AP sempre para em no máximo $n + 1$ passos.
    \item \textbf{Terminação do ALL ($O(n^2)$):} cada iteração do ciclo principal ($q_1/q_8 \to q_2/q_3 \to q_4/q_5 \to q_6 \to q_7 \to q_8$) converte irreversivelmente um par de $\{\texttt{i}, \texttt{s}\}$ (um de cada lado de \texttt{:}) em $\{\texttt{I}, \texttt{S}\}$, sem nunca desmarcar células. Entre marcações, cada varredura move o cabeçote em direção monotônica sobre no máximo $n + 1$ células. Com no máximo $\lfloor n/2 \rfloor$ iterações, o ALL sempre para incondicionalmente em no máximo $(n+2)^2$ passos.
\end{itemize}
